\documentclass[12pt, a4paper]{article}

% ── Packages ──────────────────────────────────────────────────────────────────
\usepackage[a4paper, margin=2.5cm]{geometry}
\usepackage{graphicx}
\usepackage{float}
\usepackage{amsmath}
\usepackage{amssymb}
\usepackage{xcolor}
\usepackage{mdframed}
\usepackage{parskip}
\usepackage{enumitem}
\usepackage{titlesec}
\usepackage{booktabs}
\usepackage{hyperref}
\usepackage[T1]{fontenc}
\usepackage{tikz}
\usepackage{listings}
\usetikzlibrary{arrows.meta}

% ── Graphics path ─────────────────────────────────────────────────────────────
\graphicspath{{../../Images_for_latex/Lecture_8/}}

% ── Colours ───────────────────────────────────────────────────────────────────
\definecolor{keyblue}{RGB}{30, 80, 160}
\definecolor{keybluebg}{RGB}{220, 235, 255}
\definecolor{noteorange}{RGB}{200, 100, 0}
\definecolor{noteorangebg}{RGB}{255, 240, 215}
\definecolor{warnred}{RGB}{180, 30, 30}
\definecolor{warnredbg}{RGB}{255, 220, 220}
\definecolor{defgreen}{RGB}{20, 110, 50}
\definecolor{defgreenbg}{RGB}{215, 245, 225}
\definecolor{headernavy}{RGB}{25, 35, 90}

% ── Box environments ──────────────────────────────────────────────────────────
\newmdenv[linecolor=keyblue,backgroundcolor=keybluebg,linewidth=1.5pt,
  leftline=true,rightline=false,topline=false,bottomline=false,
  innerleftmargin=10pt,innerrightmargin=8pt,
  innertopmargin=6pt,innerbottommargin=6pt,
  skipabove=8pt,skipbelow=8pt]{keybox}

\newmdenv[linecolor=noteorange,backgroundcolor=noteorangebg,linewidth=1.5pt,
  leftline=true,rightline=false,topline=false,bottomline=false,
  innerleftmargin=10pt,innerrightmargin=8pt,
  innertopmargin=6pt,innerbottommargin=6pt,
  skipabove=8pt,skipbelow=8pt]{notebox}

\newmdenv[linecolor=warnred,backgroundcolor=warnredbg,linewidth=1.5pt,
  leftline=true,rightline=false,topline=false,bottomline=false,
  innerleftmargin=10pt,innerrightmargin=8pt,
  innertopmargin=6pt,innerbottommargin=6pt,
  skipabove=8pt,skipbelow=8pt]{warnbox}

\newmdenv[linecolor=defgreen,backgroundcolor=defgreenbg,linewidth=1.5pt,
  leftline=true,rightline=false,topline=false,bottomline=false,
  innerleftmargin=10pt,innerrightmargin=8pt,
  innertopmargin=6pt,innerbottommargin=6pt,
  skipabove=8pt,skipbelow=8pt]{defbox}

% ── Section formatting ────────────────────────────────────────────────────────
\titleformat{\section}[block]{\large\bfseries\color{headernavy}}{}{0pt}{}
  [\vspace{2pt}\rule{\linewidth}{0.8pt}\vspace{4pt}]
\titleformat{\subsection}[block]{\normalsize\bfseries\color{keyblue}}{}{0pt}{}

% ── Document ──────────────────────────────────────────────────────────────────
\begin{document}

\begin{center}
  {\LARGE\bfseries\color{headernavy} Sensors \& Systems --- Lecture 8 Notes}\\[4pt]
  {\large Trilateration-based Sensors: BLE, Games-on-Track, GPS}\\[2pt]
  {\normalsize Jan Bendtsen, Torben Knudsen \quad|\quad Electronic Systems, 8th Semester}
\end{center}

\vspace{4pt}
\hrule height 1.2pt
\vspace{12pt}

% ══════════════════════════════════════════════════════════════════════════════
\section*{Trilateration}
% ══════════════════════════════════════════════════════════════════════════════

\begin{figure}[H]
  \centering
  \includegraphics[width=0.65\textwidth]{slide-03.jpg}
  \caption{Three anchors at known positions each measure a range $r_i$
    to the unknown target via a Poll/Response exchange.  Each range defines
    a circle; the intersection of all three gives the unique target location.}
\end{figure}

\textbf{Trilateration} estimates an unknown position from distance
measurements to \textbf{anchors} at known coordinates.
One range places the target on a circle (2D) or sphere (3D).
A second narrows this to two candidate points.
A \textbf{third} gives the unique solution.
In 3D, four anchors are needed when an additional unknown exists
(e.g.\ the receiver clock error in GPS).

\clearpage
% ══════════════════════════════════════════════════════════════════════════════
\section*{Underwater Acoustic Positioning}
% ══════════════════════════════════════════════════════════════════════════════

\begin{figure}[H]
  \centering
  \includegraphics[width=0.72\textwidth]{slide-07.jpg}
  \caption{Underwater sensor network.  Acoustic signals handle ranging
    between underwater nodes; radio relays the computed position to shore.}
\end{figure}

The same trilateration principle applies underwater, but radio is not
an option --- salt water absorbs radio waves within a few centimetres.
\textbf{Acoustic signals} (sound) are used instead, with $c \approx
1500$\,m/s, giving the same $r = c\,\Delta t / 2$.

\begin{notebox}
  \textbf{Two media, two roles --- neither can replace the other.}
  \textbf{Acoustic} does the ranging underwater (sound travels hundreds
  of metres; radio absorbed within cm of salt water).
  \textbf{Radio} relays the already-computed position from a surface
  buoy to a satellite or shore station (sound too slow and short-range
  in air; radio travels at $c$ over thousands of km).
  The surface buoy is the bridge --- listening acoustically below,
  transmitting by radio above.
\end{notebox}

% ══════════════════════════════════════════════════════════════════════════════
\section*{GPS --- Coordinate Systems}
% ══════════════════════════════════════════════════════════════════════════════

\subsection*{ECEF and local coordinates}

\begin{figure}[H]
  \centering
  \includegraphics[width=0.55\textwidth]{slide-09.jpg}
  \caption{ECEF coordinate system.  Origin at Earth's centre; $z$
    toward the north pole; $x$ toward the Greenwich meridian.
    The frame rotates with the Earth so fixed surface points always
    have the same $(x,y,z)$.}
\end{figure}

All GPS calculations use \textbf{ECEF} (Earth-Centred Earth-Fixed) ---
a Cartesian system rotating with the Earth.
A point at longitude $\lambda$, latitude $\phi$, radial distance
$\gamma$ from Earth's centre has ECEF coordinates:
\[
  \begin{bmatrix} x \\ y \\ z \end{bmatrix}
  = \gamma
    \begin{bmatrix}
      \cos\phi\cos\lambda \\
      \cos\phi\sin\lambda \\
      \sin\phi
    \end{bmatrix}
\]
where $\lambda$ is the angle east/west of Greenwich ($-180^\circ$ to
$+180^\circ$), $\phi$ is the angle north/south of the equator
($-90^\circ$ to $+90^\circ$), and $\gamma \approx R_E = 6{,}371$\,km
for points on the surface.

ECEF is good for the maths but hard to interpret in practice.
\textbf{Local frames} centred on the observer make errors intuitive:

\begin{itemize}[leftmargin=1.2em,itemsep=2pt]
  \item \textbf{NED} (North-East-Down): standard in aviation; $z$
    points into the Earth so altitude increases as $z$ decreases.
  \item \textbf{ENU} (East-North-Up): $z$ points up, more natural
    for ground use.
  \item \textbf{AER} (Azimuth-Elevation-Range): describes a target
    by compass direction, elevation angle above horizon, and slant
    distance --- used for pointing at satellites.
\end{itemize}

\begin{figure}[H]
  \centering
  \includegraphics[width=0.82\textwidth]{slide-12.jpg}
  \caption{NED, ENU, and AER local coordinate systems, all centred at
    a reference point on the surface.}
\end{figure}

% ══════════════════════════════════════════════════════════════════════════════
\section*{GPS --- Pseudoranges and Why 4 Satellites}
% ══════════════════════════════════════════════════════════════════════════════

\begin{figure}[H]
  \centering
  \includegraphics[width=0.65\textwidth]{slide-08.jpg}
  \caption{GPS trilateration in 3D.  Three spheres give two candidate
    points; a fourth satellite resolves this and simultaneously solves
    for the unknown receiver clock error.}
\end{figure}

Standard 3D trilateration needs 3 ranges for 3 unknowns $(x,y,z)$.
GPS needs a \textbf{fourth satellite} because the receiver clock error
is a fourth unknown.
The satellite carries an \textbf{atomic clock} ($< 1$\,\textmu s/day
drift); the receiver uses a cheap \textbf{quartz oscillator} (tens of
\textmu s/second drift).
Because the receiver clock is off by unknown $\tau$, every range
measurement is corrupted by the same $c\tau$, giving a
\textbf{pseudorange}:
\[
  \boxed{P_S = \rho(t,\,t_S) + c\tau - c\tau_S}
\]

\begin{center}
\renewcommand{\arraystretch}{1.4}
\begin{tabular}{p{1.3cm} p{3cm} p{7.5cm}}
  \toprule
  \textbf{Symbol} & \textbf{Name} & \textbf{Meaning} \\
  \midrule
  $P_S$           & Pseudorange [m]       & Measured apparent distance: raw $\Delta t \times c$ \\
  $\rho(t,\,t_S)$ & True geometric range  & $\sqrt{(x_S{-}x)^2 + (y_S{-}y)^2 + (z_S{-}z)^2}$ \\
  $\tau$          & Receiver clock error  & Unknown offset of receiver quartz clock from GPS time \\
  $\tau_S$        & Satellite clock error & Known --- provided in Navigation Message as polynomial \\
  $t,\,t_S$       & True times            & True GPS times of reception and transmission \\
  \bottomrule
\end{tabular}
\end{center}

\begin{notebox}
  \textbf{Clock drift and the Navigation Message.}
  Even the satellite's atomic clock drifts slightly.
  The GPS Control Segment (a global network of ground stations) tracks
  all satellites, measures their clock deviations, and uplinks
  correction coefficients packed into the Navigation Message as:
  \[
    \delta t^s = a_0 + a_1(t - t_0) + a_2(t - t_0)^2
  \]
  where $a_0, a_1, a_2$ are the offset, drift, and drift-rate
  coefficients, and $t_0$ is the reference epoch.
  The receiver reads and subtracts this correction, so satellite clock
  error is essentially a solved problem.
  The receiver clock error $\tau$ has no such correction and must be
  solved for alongside $(x,y,z)$ --- hence 4 satellites.
\end{notebox}

% ══════════════════════════════════════════════════════════════════════════════
\section*{GPS --- Signal and Cross-Correlation}
% ══════════════════════════════════════════════════════════════════════════════

\begin{figure}[H]
  \centering
  \includegraphics[width=0.72\textwidth]{slide-23.jpg}
  \caption{GPS signal structure.  The L1 carrier (1575.42\,MHz) is
    phase-modulated by the XOR of the C/A ranging code (1\,Mbps) and
    the Navigation Message (50\,bps).}
\end{figure}

The GPS signal is built in three layers, all derived from the atomic
clock at $f_0 = 10.23$\,MHz:

\begin{itemize}[leftmargin=1.2em,itemsep=3pt]
  \item \textbf{L1 carrier} ($154 \times f_0 = 1575.42$\,MHz): the
    radio wave that carries everything through 20\,000\,km of space.
  \item \textbf{C/A code} ($f_0/10 = 1.023$\,Mbps): a 1023-chip
    pseudorandom sequence repeating every millisecond.
    Each satellite uses a \emph{different} code (CDMA), allowing all
    satellites to share the same frequency without interfering.
    The fast chip rate gives microsecond-precision timing for ranging.
  \item \textbf{Navigation Message} (50\,bps): the slow data payload
    carrying clock corrections, ephemeris, and almanac.
    Timing comes from the C/A code, not from these bits.
\end{itemize}

\begin{figure}[H]
  \centering
  \includegraphics[width=0.72\textwidth]{slide-24.jpg}
  \caption{Signal reception via cross-correlation.  The receiver
    generates a local replica of the satellite's PRN code and slides
    it in time.  When misaligned the correlation is $\approx 0$;
    when aligned it spikes to $+1$, giving the clock difference
    $T - T^s$ and hence the pseudorange.}
\end{figure}

The receiver measures pseudorange via \textbf{cross-correlation}:
\[
  R(k) = \frac{1}{N}\sum_{i=0}^{N-1} x_i \cdot y_{i-k}
\]
where $x_i$ is the received chip, $y_{i-k}$ is the local replica
delayed by $k$ chips, and $N$ is the number of chips summed.
When the lag $k \neq \Delta t$ the pseudorandom sequences are
misaligned and products nearly cancel ($R \approx 0$).
When $k = \Delta t$ every chip multiplies its own copy giving
$R = +1$ --- a sharp spike.
The lag at the spike is the travel time:
\[
  \Delta t = k \qquad \Rightarrow \qquad \rho = c\,\Delta t
\]
This is exactly the same principle as Games-on-Track (chirp instead
of PRN code): transmit a known pattern, generate a local copy, slide
until they align --- the required shift is the travel time.

% ══════════════════════════════════════════════════════════════════════════════
\section*{GPS --- Position Estimation and Error}
% ══════════════════════════════════════════════════════════════════════════════

\subsection*{Linearisation and least squares}

\begin{figure}[H]
  \centering
  \includegraphics[width=0.72\textwidth]{slide-26.jpg}
  \caption{With $m \geq 4$ satellites, the pseudorange equations are
    linearised around an initial guess giving $b = A\hat{\xi} + \nu$.}
\end{figure}

With $m$ satellites and satellite clock corrections applied, the
pseudorange equations become:
\[
  \tilde{P}_k = \rho_k(x,y,z) + c\tau, \qquad k = 1,\dots,m
\]
This is nonlinear in the unknowns $(x,y,z,\tau)$.
Linearising around an initial guess $(x_0,y_0,z_0,\tau_0)$ and
defining the correction vector $\hat{\xi}$ and measurement residual
$b$:
\[
  \hat{\xi} =
  \begin{bmatrix} x - x_0 \\ y - y_0 \\ z - z_0 \\ \tau - \tau_0
  \end{bmatrix}, \qquad
  b =
  \begin{bmatrix} \tilde{P}_1 - P_{0,1} \\
    \vdots \\ \tilde{P}_m - P_{0,m}
  \end{bmatrix}
\]
gives the linear system $b = A\hat{\xi} + \nu$, where $\nu$ is the
fit residual.
Each row of the Jacobian $A$ is a unit vector pointing from the
receiver toward satellite $k$, with $c$ appended for the clock
column:
\[
  A_k =
  \begin{bmatrix}
    -\dfrac{x_S^{(k)}-x_0}{\rho_{0,k}} &
    -\dfrac{y_S^{(k)}-y_0}{\rho_{0,k}} &
    -\dfrac{z_S^{(k)}-z_0}{\rho_{0,k}} & c
  \end{bmatrix}
\]
where $\rho_{0,k}$ is the predicted range to satellite $k$ from the
initial guess, and $(x_S^{(k)}, y_S^{(k)}, z_S^{(k)})$ is the known
ECEF position of satellite $k$.
With $m > 4$ the system is overdetermined; minimising
$\|\nu\|^2$ gives the \textbf{least-squares solution}:
\[
  \boxed{\hat{\xi} = (A^\top A)^{-1} A^\top b}
\]

\begin{notebox}
  \textbf{Why more satellites help and why geometry matters (DOP).}
  With exactly 4 satellites any noise goes straight into the position
  estimate.
  With $m > 4$, extra measurements are redundant but useful --- least
  squares averages out random noise across all satellites.
  The conditioning of $A^\top A$ quantifies satellite geometry:
  good spread across the sky $\rightarrow$ $A^\top A$ well-conditioned
  $\rightarrow$ small position errors.
  Poor spread $\rightarrow$ $A^\top A$ nearly singular $\rightarrow$
  large errors.
  This is exactly what \textbf{DOP} (Dilution of Precision)
  measures --- a high DOP means bad geometry and amplified errors.
\end{notebox}

\subsection*{Error propagation and covariance}

\begin{figure}[H]
  \centering
  \includegraphics[width=0.68\textwidth]{slide-28.jpg}
  \caption{Pseudorange measurement noise $\Sigma_\nu$ propagated
    through the least-squares formula gives the position covariance
    $\Sigma_\xi$.}
\end{figure}

The estimation error propagates from pseudorange noise $\nu$ as:
\[
  \nu_\xi = (A^\top A)^{-1} A^\top \nu
\]
Taking the expectation of the outer product gives the
\textbf{covariance matrix} of position errors:
\[
  \Sigma_\xi = (A^\top A)^{-1} A^\top \Sigma_\nu A\,(A^\top A)^{-1}
\]
where $\Sigma_\nu$ is the $m \times m$ covariance of the pseudorange
noise.
Assuming equal, independent noise on all satellites
($\Sigma_\nu = \sigma_\nu^2 I$) this simplifies to:
\[
  \Sigma_\xi = \sigma_\nu^2\,(A^\top A)^{-1}
\]
The \textbf{diagonal elements} give individual variances:
$\sigma_x^2, \sigma_y^2, \sigma_z^2$ for position and $\sigma_\tau^2$
for the clock error estimate.
The shape of the uncertainty ellipsoid is determined entirely by
satellite geometry $(A^\top A)^{-1}$; noise level $\sigma_\nu^2$
scales the whole thing uniformly.

\subsection*{Uncertainty ellipsoid in local coordinates}

\begin{figure}[H]
  \centering
  \includegraphics[width=0.65\textwidth]{slide-29.jpg}
  \caption{The ECEF covariance $\Sigma_\xi$ is rotated into local
    North-East-Up coordinates using $R_{\phi\lambda}$ for a more
    intuitive decomposition into horizontal and vertical errors.}
\end{figure}

An error of 5\,m in the ECEF $x$-direction is hard to interpret.
Rotating into local North-East-Up (NEU) coordinates makes it
meaningful:
\[
  \nu_\text{NEU} = R_{\phi\lambda}\,\nu_\xi, \qquad
  \Sigma_\text{NEU} = R_{\phi\lambda}\,\Sigma_\xi\,R_{\phi\lambda}^\top
\]
where the rotation matrix at latitude $\phi$ and longitude $\lambda$
is:
\[
  R_{\phi\lambda} =
  \begin{bmatrix}
    -\sin\phi\cos\lambda & -\sin\phi\sin\lambda &  \cos\phi \\
    -\sin\lambda         &  \cos\lambda         &  0        \\
     \cos\phi\cos\lambda &  \cos\phi\sin\lambda &  \sin\phi
  \end{bmatrix}
\]
The diagonal of $\Sigma_\text{NEU}$ now gives
$\sigma_\text{North}^2$, $\sigma_\text{East}^2$,
$\sigma_\text{Up}^2$ --- directly telling you the horizontal spread
and vertical offset of the uncertainty ellipsoid.

\subsection*{Tracking moving targets with a Kalman filter}

\begin{figure}[H]
  \centering
  \includegraphics[width=0.78\textwidth]{slide-30.jpg}
  \caption{A Kalman filter takes consecutive GPS position estimates
    and infers both position and velocity using a constant-velocity
    motion model.}
\end{figure}

GPS gives a new position estimate every second; a
\textbf{Kalman filter} turns this into a smooth position-and-velocity
track.
With no other information, a constant-velocity model is used:
\[
  \chi_k =
  \begin{bmatrix} x_k \\ y_k \\ z_k \\ v_{x,k} \\ v_{y,k} \\ v_{z,k}
  \end{bmatrix}, \quad
  \Phi =
  \begin{bmatrix}
    1 & 0 & 0 & T_s & 0   & 0   \\
    0 & 1 & 0 & 0   & T_s & 0   \\
    0 & 0 & 1 & 0   & 0   & T_s \\
    0 & 0 & 0 & 1   & 0   & 0   \\
    0 & 0 & 0 & 0   & 1   & 0   \\
    0 & 0 & 0 & 0   & 0   & 1
  \end{bmatrix}, \quad
  H = \begin{bmatrix} I & 0 \end{bmatrix} \in \mathbb{R}^{3\times6}
\]
where $\chi_k$ is the state vector of position and velocity,
$\Phi$ is the state transition matrix with sample time $T_s$,
and $H$ extracts position only (GPS measures position, not velocity
directly --- the filter infers velocity from how position changes
between steps).
The covariance $\Sigma_\nu$ from the previous section is used directly
as the measurement noise covariance in the Kalman filter.

\begin{keybox}
  \textbf{Lecture 8 --- Key Takeaways}
  \begin{itemize}[leftmargin=1.2em,itemsep=2pt]
    \item \textbf{BLE RSSI}: range via $\hat{\rho} = 10^{(A-S)/(10\gamma)}$.
      No timing hardware; 1--5\,m accuracy indoors due to multipath
      and $\gamma$ uncertainty.
    \item \textbf{Games-on-Track / GPS}: TOF via cross-correlation of a
      known pattern (chirp / PRN code). Locks onto first-arriving signal;
      robust to multipath; needs dedicated timing hardware.
    \item \textbf{Underwater}: radio replaced by acoustic
      ($c \approx 1500$\,m/s); same TOF principle. Radio used only at
      surface to relay position to shore.
    \item \textbf{GPS pseudorange}: $P_S = \rho + c\tau - c\tau_S$.
      Satellite clock error $\tau_S$ corrected via Navigation Message
      polynomial. Receiver clock error $\tau$ solved alongside
      $(x,y,z)$ --- requires 4 satellites minimum.
    \item \textbf{Least squares}: $\hat{\xi} = (A^\top A)^{-1}A^\top b$,
      where each row of $A$ is a unit vector toward a satellite with $c$
      appended. More satellites reduce noise; satellite geometry (DOP)
      determines how much noise gets amplified into position error.
    \item \textbf{Covariance}: $\Sigma_\xi = \sigma_\nu^2(A^\top A)^{-1}$;
      rotated to local NEU for interpretable horizontal/vertical errors.
      Fed directly into a Kalman filter for tracking moving targets.
    \item \textbf{ECEF} rotates with Earth; local frames (NED, ENU, AER)
      express errors relative to the observer.
  \end{itemize}
\end{keybox}

\end{document}
